\documentclass[11pt,a4paper]{article}
\usepackage{fontspec}
\usepackage[sfdefault]{FiraSans}
\usepackage[a4paper,margin=1.8cm,top=2.4cm,bottom=2cm]{geometry}
\usepackage{xcolor}
\definecolor{nvgreen}{HTML}{76B900}
\definecolor{nvdgreen}{HTML}{4F7A00}
\definecolor{nvdark}{HTML}{1A1A2E}
\definecolor{nvcard}{HTML}{2D2D44}
\definecolor{nvlightgreen}{HTML}{A3D944}
\definecolor{nvgray}{HTML}{6B6B80}
\definecolor{nvorange}{HTML}{D35400}
\definecolor{nvred}{HTML}{C0392B}
\usepackage{listings}
\usepackage[most]{tcolorbox}
\usepackage{booktabs}
\usepackage{colortbl}
\usepackage{tabularx}
\usepackage{amssymb}
\usepackage{titlesec}
\usepackage{fancyhdr}
\usepackage{enumitem}
\usepackage{hyperref}
\hypersetup{colorlinks=true, urlcolor=nvdgreen, linkcolor=nvdgreen}
\setlist{leftmargin=1.4em, topsep=2pt, itemsep=1pt}

\titleformat{\section}{\Large\bfseries\color{nvgreen}}{\thesection}{0.5em}{}[{\color{nvgreen}\titlerule[1.2pt]}]
\titlespacing{\section}{0pt}{10pt}{5pt}
\titleformat{\subsection}{\large\bfseries\color{nvdgreen}}{\thesubsection}{0.4em}{}
\titlespacing{\subsection}{0pt}{7pt}{3pt}

\lstdefinestyle{nv}{
  backgroundcolor=\color{nvdark}, basicstyle=\ttfamily\footnotesize\color{white},
  keywordstyle=\color{nvlightgreen}, commentstyle=\color{nvgray}\itshape,
  stringstyle=\color{nvlightgreen}, breaklines=true, frame=single, framerule=0pt,
  xleftmargin=7pt, xrightmargin=4pt, aboveskip=6pt, belowskip=6pt,
  showstringspaces=false, columns=fullflexible, keepspaces=true, breakindent=0pt,
}
\lstset{style=nv}

\newtcolorbox{warn}{colback=nvorange!7,colframe=nvorange,boxrule=1pt,arc=2pt,left=7pt,right=7pt,top=4pt,bottom=4pt,fontupper=\small}
\newtcolorbox{tip}{colback=nvgreen!8,colframe=nvgreen,boxrule=1pt,arc=2pt,left=7pt,right=7pt,top=4pt,bottom=4pt,fontupper=\small}

\pagestyle{fancy}\fancyhf{}
\fancyhead[L]{\small\color{nvgray}Runbook Deployment SLM Spesialis}
\fancyhead[R]{\small\color{nvgray}Module 04 $\cdot$ Navasena}
\fancyfoot[C]{\small\color{nvgray}\thepage}
\renewcommand{\headrulewidth}{0.4pt}\renewcommand{\headrule}{\hbox to\headwidth{\color{nvgreen}\leaders\hrule height \headrulewidth\hfill}}

\newcommand{\code}[1]{\colorbox{nvbg}{\texttt{\small #1}}}
\definecolor{nvbg}{HTML}{EFF3E4}

\begin{document}
\thispagestyle{fancy}

\noindent\begin{tcolorbox}[colback=nvgreen,colframe=nvgreen,arc=4pt,left=12pt,right=12pt,top=8pt,bottom=8pt]
\color{white}\LARGE\bfseries Runbook Deployment SLM Spesialis\\[3pt]
\normalsize\mdseries Mac + Jetson Orin Nano 8GB (Edge-First) $\cdot$ Module 04 Capstone $\cdot$ Navasena Gen-ML Course
\end{tcolorbox}

\vspace{4pt}
\noindent Panduan langkah-demi-langkah men-deploy SLM spesialis (hasil notebook 05/06/07) ke perangkat nyata. Semua perintah \& angka di sini \textbf{diverifikasi langsung} di MacBook M4 dan Jetson Orin Nano 8GB (2026-06-12).

\section*{\color{nvgreen}\S0\quad Prasyarat \& Artefak}
\textbf{Artefak} (hasil notebook 07, sudah di HuggingFace):
\begin{itemize}
  \item Repo GGUF \texttt{chmdznr/qwen3-0.6b-spesialis-gguf} (Q8\_0 $\sim$0.6 GB, Q4\_K\_M $\sim$0.4 GB) --- \textbf{publik}, device mana pun bisa tarik tanpa login.
  \item \texttt{Modelfile} (\S C) --- berisi TEMPLATE ChatML + SYSTEM prompt.
\end{itemize}
\textbf{Perangkat (diverifikasi):} MacBook Apple Silicon (M4, macOS 26) backend \textbf{Metal}; Jetson Orin Nano 8GB Super, JetPack 6.2.1 / CUDA 12.6, $\sim$5.6 GB RAM bebas, backend \textbf{CUDA}.

\textbf{Prinsip jembatan:} HuggingFace Hub jadi distribusi. Push GGUF sekali dari Colab $\to$ tiap device tinggal \texttt{ollama create} yang menarik dari HF. \textbf{Nol scp, nol USB.}

\begin{warn}
\textbf{Tiga jebakan yang sudah dibuktikan} (baca dulu):
\begin{enumerate}[leftmargin=1.6em]
  \item \texttt{ollama run hf.co/...} telanjang \textbf{KEHILANGAN system prompt} $\to$ spesialis balik jadi chatbot biasa. \textbf{Wajib \texttt{ollama create} dengan Modelfile ber-SYSTEM.}
  \item \textbf{Ollama \textless\ 0.30.7 menyisakan \texttt{<think></think>} kosong} (thinking-token Qwen3). \textbf{Upgrade ke $\geq$ 0.30.7} $\to$ bersih.
  \item \textbf{Pakai Q8\_0} untuk spesialis mungil ($\sim$0.6 GB) --- near-lossless \& aman dari bug memori L4T.
\end{enumerate}
\end{warn}

\section*{\color{nvgreen}\S A\quad MacBook (Apple Silicon) --- dev lokal}
\textbf{A1. Install Ollama}
\begin{lstlisting}
brew install ollama        # atau app dari ollama.com/download
ollama serve &             # backend Metal otomatis
ollama --version           # pastikan >= 0.30.7
\end{lstlisting}
\textbf{A2--A3. Create dari HF (Modelfile \S C) + uji}
\begin{lstlisting}
ollama create navasena-spesialis -f Modelfile   # FROM menarik GGUF dari HF (~0.6 GB)
ollama run navasena-spesialis "Pesan 3 kopi, harga 25000 per item."
# -> {"produk":"kopi","jumlah":3,"harga_satuan":25000,"total":75000}
\end{lstlisting}
\textbf{A4. Benchmark}
\begin{lstlisting}
ollama run navasena-spesialis --verbose "minta 7 beras 14000 satunya"  # baca "eval rate"
ollama ps    # PROCESSOR harus "100% GPU"
\end{lstlisting}
\begin{tip}\textbf{Hasil terverifikasi Mac M4:} 3/3 ekstraksi JSON benar, \textbf{$\sim$85 tok/dtk}, output bersih (Ollama 0.30.7).\end{tip}

\section*{\color{nvgreen}\S B\quad Jetson Orin Nano 8GB --- target edge}
\textbf{B1. Masuk + cek}
\begin{lstlisting}
ssh jetson@<jetson-ip>      # mis. 192.168.1.114
ollama --version ; free -h
\end{lstlisting}
\textbf{B2. Upgrade Ollama ke $\geq$ 0.30.7} (kalau masih 0.30.6 $\to$ output kotor \texttt{<think>})
\begin{lstlisting}
# Jalankan SELURUH script sebagai root agar tak ada prompt sudo internal (butuh tty):
sudo bash -c "curl -fsSL https://ollama.com/install.sh | sh"
# Installer otomatis ambil build khusus: ollama-linux-arm64-jetpack6.tar.zst (butuh paket zstd)
ollama --version            # -> 0.30.7+
\end{lstlisting}
\textbf{B3--B4. Create dari HF + verifikasi GPU}
\begin{lstlisting}
mkdir -p ~/navasena && cd ~/navasena && nano Modelfile   # tempel isi \S C
ollama create navasena-spesialis -f Modelfile
ollama run navasena-spesialis "Tolong kirim 12 sabun harga 8500 per item."
# -> {"produk":"sabun","jumlah":12,"harga_satuan":8500,"total":102000}
ollama ps    # PROCESSOR -> "100% GPU" (CUDA)
\end{lstlisting}
\textbf{B5. Termal/daya + benchmark}
\begin{lstlisting}
sudo nvpmodel -m 0          # MAXN (atau -m 1 utk 25W)
sudo jetson_clocks ; sudo tegrastats     # pantau RAM/GPU%/suhu (Ctrl-C berhenti)
ollama run navasena-spesialis --verbose "minta 7 beras 14000 satunya"
\end{lstlisting}
\begin{tip}\textbf{Hasil terverifikasi Jetson Orin Nano:} 3/3 JSON benar, \textbf{$\sim$39 tok/dtk}, 100\% GPU, output bersih setelah upgrade 0.30.7.\end{tip}
\textbf{B6. (Opsional) Ekspose ke LAN} agar Mac bisa panggil Jetson
\begin{lstlisting}
sudo systemctl edit ollama   # tambah: [Service]  Environment="OLLAMA_HOST=0.0.0.0:11434"
sudo systemctl restart ollama
\end{lstlisting}
\textbf{B7. (Advanced) Build llama.cpp-CUDA}
\begin{lstlisting}
which nvcc && nvcc --version   # WAJIB 12.6 di PATH
git clone https://github.com/ggml-org/llama.cpp && cd llama.cpp
cmake -B build -DGGML_CUDA=ON -DCMAKE_CUDA_ARCHITECTURES="87"   # sm_87 = Orin
cmake --build build --config Release --parallel 4
\end{lstlisting}

\section*{\color{nvgreen}\S C\quad Modelfile (identik Mac \& Jetson)}
\begin{warn}Ollama \textbf{mengabaikan} chat-template Jinja di GGUF $\to$ kita suplai TEMPLATE ChatML sendiri. \textbf{SYSTEM wajib} --- tanpa itu spesialis balik jadi chatbot biasa.\end{warn}
\begin{lstlisting}
FROM hf.co/chmdznr/qwen3-0.6b-spesialis-gguf:Q8_0

TEMPLATE """{{- if .System }}<|im_start|>system
{{ .System }}<|im_end|>
{{ end }}{{- range .Messages }}<|im_start|>{{ .Role }}
{{ .Content }}<|im_end|>
{{ end }}<|im_start|>assistant
"""

PARAMETER stop "<|im_start|>"
PARAMETER stop "<|im_end|>"
PARAMETER temperature 0.2

SYSTEM Ekstrak pesanan ke JSON dengan key produk, jumlah, harga_satuan, total; angka tanpa titik; keluarkan HANYA JSON.
\end{lstlisting}

\section*{\color{nvgreen}\S D\quad Pola Client Universal (OpenAI-compatible)}
Inti deployment: \textbf{kode aplikasi portabel.} Ollama (lokal/edge), vLLM (cloud), NVIDIA NIM --- semua OpenAI-compatible. Pindah target = ganti \textbf{3 baris}.
\begin{lstlisting}
from openai import OpenAI
client = OpenAI(base_url="http://localhost:11434/v1/", api_key="ollama")  # Ollama lokal (Mac)
# Jetson di LAN :  base_url="http://<jetson-ip>:11434/v1"
# vLLM cloud    :  base_url="http://<host>:8000/v1", api_key="EMPTY"
# NVIDIA NIM    :  base_url="https://integrate.api.nvidia.com/v1", api_key=<nvapi-key>
MODEL = "navasena-spesialis"
r = client.chat.completions.create(model=MODEL,
    messages=[{"role":"user","content":"Pesan 5 kopi, harga 25000 per item."}])
print(r.choices[0].message.content)
\end{lstlisting}

\section*{\color{nvgreen}\S E\quad Perbandingan \& Kerangka Keputusan}
\renewcommand{\arraystretch}{1.25}
\noindent\begin{tabularx}{\textwidth}{@{}l l l c X@{}}
\toprule
\rowcolor{nvgreen}\color{white}\textbf{Tool} & \color{white}\textbf{Target} & \color{white}\textbf{Setup} & \color{white}\textbf{Course?} & \color{white}\textbf{Kapan dipakai} \\
\midrule
\textbf{Ollama} & Mac, Jetson & termudah & \textcolor{nvgreen}{\checkmark} & default edge/dev \\
\rowcolor{nvbg}\textbf{llama.cpp} & Mac, Jetson & build manual & \textcolor{nvgreen}{\checkmark} opsi & kontrol penuh, benchmark mentah \\
TensorRT-LLM & Jetson produksi & build engine berat & konsep & latency-kritis fixed-model \\
\rowcolor{nvbg}vLLM / SGLang & GPU datacenter & server & konsep & banyak user, throughput \\
NVIDIA NIM & cloud & API & \textcolor{nvgreen}{\checkmark} & model besar via API \\
\rowcolor{nvbg}TGI & --- & --- & legacy & \emph{maintenance mode sejak Des 2025} \\
\bottomrule
\end{tabularx}

\vspace{4pt}\noindent\textbf{Pohon keputusan:} 1 user/edge $\to$ Ollama; banyak/cloud $\to$ vLLM/NIM. Target laptop/Jetson $\to$ Ollama (mudah) atau llama.cpp (dalam). \textbf{Punchline:} semua bicara OpenAI API $\to$ kode aplikasi sama, beda hanya \texttt{base\_url}.
\begin{warn}\textbf{Jujur:} vLLM/TGI \textbf{bukan} untuk Orin Nano 8GB (tool datacenter). Untuk edge: Ollama/llama.cpp/TensorRT.\end{warn}

\section*{\color{nvgreen}\S F\quad Troubleshooting}
\renewcommand{\arraystretch}{1.3}
\noindent\begin{tabularx}{\textwidth}{@{}X X X@{}}
\toprule
\rowcolor{nvgreen}\color{white}\textbf{Gejala} & \color{white}\textbf{Sebab} & \color{white}\textbf{Fix} \\
\midrule
Jawaban percakapan, bukan JSON & \texttt{ollama run hf.co/...} telanjang $\to$ SYSTEM hilang & \texttt{ollama create} dgn Modelfile ber-SYSTEM (\S C) \\
\rowcolor{nvbg}Output \texttt{<think></think>\{json\}} & Ollama \textless\ 0.30.7 tak parse thinking-token & upgrade $\geq$ 0.30.7; atau parser JSON klien (regex) tetap valid \\
Output \texttt{GGGG...} / gibberish & merge adapter ke base 4-bit (NaN) & merge dari base fp16 lalu re-export GGUF \\
\rowcolor{nvbg}\texttt{sudo: a terminal is required} & cache sudo tak dihormati installer & \texttt{sudo bash -c "curl ... | sh"} (script sbg root) \\
Installer minta \texttt{zstd} & paket zstd hilang & \texttt{sudo apt-get install -y zstd} \\
\rowcolor{nvbg}\texttt{NvMap ... ENOMEM} saat load & bug CMA/NvMap L4T 36.4 ($>$1.3 GB) & GGUF Q4 \textless\ 1.3 GB; satu model/sesi (0.6B Q8\_0 aman) \\
\texttt{ollama ps} tampil CPU & model $>$ VRAM / runtime salah & cek \texttt{nvpmodel}/\texttt{tegrastats}; pakai Q4\_K\_M \\
\rowcolor{nvbg}Build llama.cpp gagal (arch) & arsitektur CUDA salah & set CUDA arch \texttt{"87"} (sm\_87) \\
Mac tak bisa akses Jetson & Ollama loopback-only & \texttt{OLLAMA\_HOST=0.0.0.0:11434} + restart (\S B6) \\
\bottomrule
\end{tabularx}

\section*{\color{nvgreen}\S G\quad Rollback / Cleanup}
\begin{lstlisting}
ollama rm navasena-spesialis
ollama rm hf.co/chmdznr/qwen3-0.6b-spesialis-gguf   # hapus cache HF (opsional)
\end{lstlisting}

\vspace{6pt}
\begin{tcolorbox}[colback=nvdark,colframe=nvgreen,arc=3pt,left=10pt,right=10pt,top=6pt,bottom=6pt]
\color{white}\textbf{\color{nvlightgreen}Yang sudah terbukti.} Satu GGUF di HF $\to$ deploy ke \textbf{Mac (Metal)} \& \textbf{Jetson (CUDA)} dengan kode \& Modelfile identik, beda hanya \texttt{base\_url}. Mac M4 $\sim$85 tok/dtk $\cdot$ Jetson Orin Nano $\sim$39 tok/dtk, keduanya 100\% GPU, 3/3 ekstraksi benar, output bersih (Ollama 0.30.7). SLM spesialis 0.6B yang dilatih gratis di Colab kini berjalan di edge $\sim$\$249 --- inilah ``kecil tapi cakap'' yang bisa diproduksikan.\\[4pt]
\textbf{Selesai. Module 04 lengkap (00--07). Selanjutnya: Module 05 --- RAG.}
\end{tcolorbox}

\end{document}
